%%% ================================================================
%%% ТИТУЛЬНЫЙ ЛИСТ ОТЧЁТА О ПРОИЗВОДСТВЕННОЙ ПРАКТИКЕ
%%% Адаптировано из SPbPU-student-thesis-template (practice.tex)
%%% Поля, отмеченные %% TODO, заполнить своими данными.
%%% ================================================================

\thispagestyle{empty}
\newgeometry{top=2cm, bottom=2cm, left=3cm, right=1cm,
             headsep=0cm, footskip=0cm}

{\centering\small
  \MakeUppercase{Федеральное государственное автономное образовательное
    учреждение высшего образования}\\[2pt]
  {\bfseries
    <<\MakeUppercase{Санкт-Петербургский политехнический университет
      Петра\,Великого}>>\\[2pt]
    \MakeUppercase{Институт компьютерных наук и~кибербезопасности}\\[2pt]
    \MakeUppercase{Высшая школа программной инженерии}
  }
\par}

\vspace{0pt plus1fill}

{\centering\bfseries
  Отчет о прохождении производственной практики\\[4pt]
  на тему: <<Разработка ML-пайплайнов обработки и анализа\\
  больших данных>>
\par}

\bigskip

{\centering\normalfont
  \uline{Михайлова Сергея Никитовича, гр.\,5130903/30302} %% TODO: проверить номер группы
\par}

\bigskip

\noindent{\bfseries Направление подготовки:}
\uline{09.03.03~Прикладная информатика}.\par\smallskip

\noindent{\bfseries Место прохождения практики:}
\uline{СПбПУ, ИКНК, ВШ~ПИ}.\par\smallskip

\noindent{\bfseries Сроки практики:}
\uline{с 11.06.2026~по 09.07.2026.}\par\smallskip %% TODO: проверить даты

\noindent{\bfseries Руководитель практики от ФГАОУ~ВО~<<СПбПУ>>:}
\uline{Пархоменко Владимир Андреевич, старший преподаватель.}\par\smallskip

\noindent{\bfseries Оценка:} \uline{\hspace*{0.15\textheight}}

\vspace{0pt plus1fill}

\noindent
\begin{tabularx}{\linewidth}{lXl}
  Руководитель практики     & & \\
  от ФГАОУ~ВО~<<СПбПУ>>    & & В.\,А.\,Пархоменко  \\[14pt]
  Обучающийся               & & С.\,Н.\,Михайлов    \\[14pt]
  Дата: \uline{09.07.2026}  & &
\end{tabularx}

\vspace{0pt plus4fill}

\restoregeometry
\newpage
